% Options for packages loaded elsewhere
\PassOptionsToPackage{unicode}{hyperref}
\PassOptionsToPackage{hyphens}{url}
\documentclass[
  11pt,
]{article}
\usepackage{xcolor}
\usepackage[margin=2.8cm]{geometry}
\usepackage{amsmath,amssymb}
\setcounter{secnumdepth}{5}
\usepackage{iftex}
\ifPDFTeX
  \usepackage[T1]{fontenc}
  \usepackage[utf8]{inputenc}
  \usepackage{textcomp} % provide euro and other symbols
\else % if luatex or xetex
  \usepackage{unicode-math} % this also loads fontspec
  \defaultfontfeatures{Scale=MatchLowercase}
  \defaultfontfeatures[\rmfamily]{Ligatures=TeX,Scale=1}
\fi
\usepackage{lmodern}
\ifPDFTeX\else
  % xetex/luatex font selection
\fi
% Use upquote if available, for straight quotes in verbatim environments
\IfFileExists{upquote.sty}{\usepackage{upquote}}{}
\IfFileExists{microtype.sty}{% use microtype if available
  \usepackage[]{microtype}
  \UseMicrotypeSet[protrusion]{basicmath} % disable protrusion for tt fonts
}{}
\makeatletter
\@ifundefined{KOMAClassName}{% if non-KOMA class
  \IfFileExists{parskip.sty}{%
    \usepackage{parskip}
  }{% else
    \setlength{\parindent}{0pt}
    \setlength{\parskip}{6pt plus 2pt minus 1pt}}
}{% if KOMA class
  \KOMAoptions{parskip=half}}
\makeatother
\usepackage{longtable,booktabs,array}
\newcounter{none} % for unnumbered tables
\usepackage{calc} % for calculating minipage widths
% Correct order of tables after \paragraph or \subparagraph
\usepackage{etoolbox}
\makeatletter
\patchcmd\longtable{\par}{\if@noskipsec\mbox{}\fi\par}{}{}
\makeatother
% Allow footnotes in longtable head/foot
\IfFileExists{footnotehyper.sty}{\usepackage{footnotehyper}}{\usepackage{footnote}}
\makesavenoteenv{longtable}
\usepackage{graphicx}
\makeatletter
\newsavebox\pandoc@box
\newcommand*\pandocbounded[1]{% scales image to fit in text height/width
  \sbox\pandoc@box{#1}%
  \Gscale@div\@tempa{\textheight}{\dimexpr\ht\pandoc@box+\dp\pandoc@box\relax}%
  \Gscale@div\@tempb{\linewidth}{\wd\pandoc@box}%
  \ifdim\@tempb\p@<\@tempa\p@\let\@tempa\@tempb\fi% select the smaller of both
  \ifdim\@tempa\p@<\p@\scalebox{\@tempa}{\usebox\pandoc@box}%
  \else\usebox{\pandoc@box}%
  \fi%
}
% Set default figure placement to htbp
\def\fps@figure{htbp}
\makeatother
\setlength{\emergencystretch}{3em} % prevent overfull lines
\providecommand{\tightlist}{%
  \setlength{\itemsep}{0pt}\setlength{\parskip}{0pt}}
\usepackage{bookmark}
\IfFileExists{xurl.sty}{\usepackage{xurl}}{} % add URL line breaks if available
\urlstyle{same}
% fallback for those not using the hyperref driver hyperxmp:
\makeatletter
\@ifundefined{xmpquote}{\newcommand{\xmpquote}[1]{#1}}{}
\makeatother
\hypersetup{
  pdftitle={Building and Evaluating a Retrieval-Augmented Generation System on FinanceBench},
  pdfauthor={Marwane Mhenni; Hamza El hamdi; Ali Cherri},
  hidelinks,
  pdfcreator={LaTeX via pandoc}}

\title{Building and Evaluating a Retrieval-Augmented Generation System
on FinanceBench}
\author{Marwane Mhenni \and Hamza El hamdi \and Ali Cherri}
\date{NLP Graded Project A}

\begin{document}
\maketitle

\section{Introduction}\label{introduction}

Retrieval-Augmented Generation (RAG) grounds a language model's answers
in a closed document collection: passages relevant to the question are
retrieved from an indexed corpus and the model is instructed to answer
from them alone. This report presents the design, implementation, and
experimental evaluation of a complete RAG pipeline for financial
question answering over SEC filings, evaluated on the FinanceBench
benchmark.

Our goal is not a single headline score but a controlled study of which
design decisions matter. We sweep 41 retrieval configurations across six
axes (chunk size, chunk overlap, document structure preservation,
embedding model, number of retrieved passages, and retrieval strategy,
including a metadata prefilter derived from the question text), evaluate
five free, locally executable generation models, and run oracle-context
and closed-book ablations that attribute end-to-end errors to retrieval
or to generation.

Three findings stand out. First, restricting the search space with
metadata parsed from the question (company and fiscal period) is the
single largest retrieval improvement we measured, worth 13 to 16 points
of recall at every k, more than any chunking or embedding change.
Second, two standard-recipe components, BM25 (alone or in hybrid fusion)
and a general-purpose cross-encoder reranker, actively hurt on this
corpus of dense financial tables. Third, the ablations show the system
is at least as generation-limited as retrieval-limited: even when handed
the gold evidence, the best local 14B model fully answers only 59
percent of the questions, and 40.7 percent of questions fail under
oracle conditions, mostly on multi-step financial arithmetic.

\section{Corpus Description}\label{corpus-description}

FinanceBench (Islam et al., 2023) is an industry benchmark for question
answering over public financial filings. Its open source split contains
150 questions over 84 documents: 10-K and 10-Q filings, 8-K reports, and
earnings releases from 40 public companies, fiscal years 2015 to 2023.
Every question carries a gold answer and one to three evidence items,
each annotated with the exact supporting text and its zero-based page
number. Questions come in three types of 50 each: metrics-generated
(read a value off a statement), domain-relevant (analyst questions
requiring judgement), and novel-generated (unusual questions unlikely to
be memorised), spanning information extraction, numerical reasoning, and
logical reasoning.

The benchmark ships a document catalog with source URLs rather than the
documents themselves, and about a quarter of those URLs have rotted
since publication. Our downloader retries with SEC-compliant request
headers and falls back to the official FinanceBench GitHub mirror; 61
documents arrived from their original URLs and 23 from the mirror, for
full coverage of all 84 documents and therefore all 150 questions. A
committed manifest records the provenance, size, and SHA-256 of every
file.

Text is extracted per page with PyMuPDF: 12,013 pages and 40.0M
characters across the 84 documents. Cleaning is deliberately
conservative (unicode NFKC normalisation, control-character removal,
whitespace bounding) because aggressive normalisation destroys the row
and column adjacency that financial tables depend on. Only 78 pages
(0.65 percent) yielded fewer than 20 characters, indicating scanned or
graphical pages; no document required OCR. Pages are the atomic unit of
ground truth throughout: retrieval is scored against the annotated
evidence pages, and a text-level containment measure catches evidence
lost inside a retrieved page.

\section{Methodology}\label{methodology}

\subsection{Segmentation}\label{segmentation}

Three chunkers cover the structure-preservation axis. The token chunker
slides fixed windows over the whole document, ignoring layout. The page
chunker treats the PDF page as the atomic unit and splits only oversized
pages. The structural chunker packs whole paragraphs greedily up to the
token budget, never cutting through one. All chunk boundaries are
measured with one fixed reference tokenizer (bert-base-uncased)
regardless of the embedding model, so that when we compare embedding
models they see literally identical chunks. Every chunk carries its
document name and the page span it covers.

\subsection{Retrieval}\label{retrieval}

Chunks are embedded with sentence-transformer models, using the query
and passage prefixes each model's authors document (E5 and BGE are
asymmetric; omitting their prefixes measurably degrades recall and would
contaminate the comparison). Normalised embeddings are indexed in FAISS
and searched by inner product. Three alternative strategies are
implemented behind the same interface: Okapi BM25; a hybrid that min-max
normalises dense and BM25 scores and combines them with weight alpha;
and cross-encoder reranking (ms-marco-MiniLM-L-6-v2) applied to a dense
candidate pool of configurable size.

The metadata filter restricts candidates to documents whose company is
named in the question, narrowed further by fiscal year when one is
mentioned. It is built by matching the catalog's 40 company names (plus
a small alias list) against the question string after both are
lowercased and stripped of punctuation; 127 of the 150 questions name
their company verbatim and the filter deactivates itself on the rest. It
reads nothing from the gold record except the question text, so it uses
only information a deployed system would have.

\subsection{Generation}\label{generation}

Answers are produced by free, locally executable instruction models
served by Ollama, prompted to answer strictly from the numbered
passages, to cite every fact as {[}n{]}, to state numeric answers with
units, and to abstain with a fixed phrase when the passages are
insufficient. The context window (16,384 tokens) sits well above the
largest measured prompt, and every response's token accounting is
checked against the window so silent truncation cannot pass unnoticed;
zero truncations were flagged across all runs. Three context modes
support the ablations: retrieved (the RAG system proper), oracle (the
gold evidence passed directly, removing retrieval from the equation),
and none (closed book, bounding how much of the score could be
pre-training memorisation).

\subsection{Evaluation}\label{evaluation}

Retrieval is scored at page level against the annotated evidence:

\begin{itemize}
\tightlist
\item
  recall@k: some retrieved chunk covers a gold evidence page;
\item
  strict recall@k: a single chunk contains at least half of an evidence
  item's 5-grams, i.e.~the evidence text itself was retrieved, not
  merely its page;
\item
  document accuracy@k, MRR@k, and nDCG@k with page-level relevance;
\item
  evidence coverage@k: the mean fraction of evidence 5-grams present
  anywhere in the retrieved set, the measure that exposes chunk-boundary
  evidence loss.
\end{itemize}

Generation is scored deterministically with scale-tolerant numeric
accuracy (FinanceBench answers are mostly dollar amounts, percentages,
and ratios; the scale-tolerant variant forgives unit mismatches such as
millions versus raw dollars but not wrong digits), SQuAD-style token F1,
citation rate and citation validity (whether a cited passage actually
covers a gold page), and abstention rate. An LLM judge (qwen2.5-14B,
temperature 0, JSON-constrained output) grades correctness (0 to 2)
against the reference answer and groundedness (0 to 2) against the
passages shown to the model. Throughout the results, accuracy means the
fraction of answers the judge scores fully correct.

Because an LLM judge is an assertion until measured, we quantify its
consistency against the deterministic numeric metric: over the 1,008
judged answers to numeric questions, the judge's fully-correct verdict
and the scale-tolerant numeric match agree on 81.6 percent of cases with
Cohen's kappa 0.607, substantial agreement given that the two measures
also disagree legitimately (an answer can state the right number for the
wrong reason, or a correct non-numeric qualification). No human
labelling was performed; we flag this in the limitations.

\section{Experimental Setup}\label{experimental-setup}

All heavy computation ran on one RunPod A100-SXM4-80GB pod (CUDA 12,
PyTorch 2.9, Python 3.12); development and the 64-test unit suite ran on
a laptop. The retrieval grid comprises 41 configurations: a 4x4 chunk
size (128, 256, 512, 1024 tokens) by overlap (0, 10, 25, 50 percent)
sweep, three structure modes, six embedding models (all-MiniLM-L6-v2,
gte-base, e5-base-v2, bge-base-en-v1.5, e5-large-v2, bge-large-en-v1.5),
k in \{1, 3, 5, 10, 20\}, BM25 at two k values, hybrid at alpha in
\{0.3, 0.5, 0.7\}, reranking with pools of 20, 50, and 100, and the
metadata filter at three k values plus one filter-plus-rerank
combination. Sweeps are one-factor-at-a-time around a fixed baseline
(token chunks of 512 with 64 overlap, bge-base, dense, k=5) so each
comparison isolates one decision. The grid completed in about 100
minutes on the A100; every run is cached under a hash of its
configuration and is resumable.

Generation consumed the top passages of the best eligible retrieval run
(the metadata-filtered dense run; its top 10 passages reach 0.627
recall, its top 20 reach 0.707) across five models spanning the four
families named in the brief and two sizes of one family: llama3.1-8B,
mistral-7B, qwen2.5-7B, phi4-14B, and qwen2.5-14B, all Q4 quantised.
Ablations add oracle and closed-book runs for both qwen sizes and a k=20
variant for qwen2.5-14B. Generation is deterministic (temperature 0),
checkpoints every ten questions, and the ten generation and judging runs
together used roughly three hours of A100 time, run three streams in
parallel against one Ollama server.

\section{Results}\label{results}

\subsection{Retrieval}\label{retrieval-1}

\begin{longtable}[]{@{}llllll@{}}
\caption{headline retrieval results (the full 41-row table is in
results/tables/retrieval\_results.md).}\tabularnewline
\toprule\noalign{}
Configuration & k & Recall@k & Strict@k & MRR@k & ms/query \\
\midrule\noalign{}
\endfirsthead
\toprule\noalign{}
Configuration & k & Recall@k & Strict@k & MRR@k & ms/query \\
\midrule\noalign{}
\endhead
\bottomrule\noalign{}
\endlastfoot
dense + filter & 20 & 0.707 & 0.507 & 0.307 & 57 \\
dense + filter & 10 & 0.627 & 0.407 & 0.301 & 37 \\
dense, plain & 20 & 0.587 & 0.427 & 0.204 & 13 \\
dense, plain & 10 & 0.500 & 0.333 & 0.198 & 8 \\
dense + filter & 5 & 0.493 & 0.300 & 0.284 & 25 \\
e5-large-v2 (best embedding) & 5 & 0.453 & 0.327 & 0.316 & 7 \\
rerank, pool 20 & 10 & 0.447 & 0.340 & 0.203 & 178 \\
baseline (bge-base) & 5 & 0.333 & 0.240 & 0.177 & 6 \\
hybrid, best alpha 0.7 & 10 & 0.320 & 0.233 & 0.152 & 341 \\
BM25 & 10 & 0.153 & 0.127 & 0.096 & 229 \\
\end{longtable}

\textbf{Metadata filtering dominates every other intervention.} At each
k the filter adds 13 to 16 recall points over plain dense search and
raises MRR by half. The mechanism is visible in the error taxonomy:
wrong-document retrievals drop to 2 percent of questions. Filtering
costs a few milliseconds and no training; on a corpus whose documents
are near-duplicates of each other in vocabulary (every 10-K contains an
income statement), identity metadata carries information that embeddings
recover only partially.

\begin{figure}
\centering
\pandocbounded{\includegraphics[keepaspectratio,alt={Recall@5 by chunk size and overlap.}]{fig_chunking.png}}
\caption{Recall@5 by chunk size and overlap.}
\end{figure}

\textbf{Chunking.} The 16-point chunk sweep spans 0.293 to 0.373 recall
at k=5, a smaller range than the filter's effect at any k. A 10 percent
overlap is best at every size; 256 and 512-token chunks lead on recall.
Evidence coverage tells the other half of the story: it doubles from
128-token chunks (0.11 to 0.15) to 1024-token chunks (0.26), because
small windows retrieve the right page but hand the generator a fragment
of the evidence. This tension (small chunks rank better, large chunks
carry more of the evidence once found) motivates the 512/64 compromise
used downstream.

\textbf{Structure preservation.} Page-bounded chunks beat token windows
(0.367 versus 0.333) with structural packing between (0.347). The gold
labels are page-level, and financial statements are laid out one per
page, so respecting the page boundary aligns chunks with both the
annotation unit and the document's own semantics.

\begin{figure}
\centering
\pandocbounded{\includegraphics[keepaspectratio,alt={Embedding models compared on identical chunks.}]{fig_embeddings.png}}
\caption{Embedding models compared on identical chunks.}
\end{figure}

\textbf{Embedding models.} e5-large-v2 is the strongest (0.453 recall
and 0.316 MRR, the best MRR of any configuration including filtered
ones), followed by bge-large (0.413), e5-base (0.387), bge-base (0.333),
gte-base (0.313), and MiniLM (0.253). The E5 family wins at equal
parameter count and scale helps within each family; yet swapping
bge-base for e5-large is still worth less than adding the metadata
filter on top of bge-base.

\begin{figure}
\centering
\pandocbounded{\includegraphics[keepaspectratio,alt={Effect of k on the four retrieval metrics.}]{fig_k.png}}
\caption{Effect of k on the four retrieval metrics.}
\end{figure}

\textbf{k.} Recall grows monotonically from 0.093 at k=1 to 0.587 at
k=20 with no plateau, so the retriever keeps finding gold pages deeper
in the ranking; the cost is prompt length and distraction for the
generator, quantified by the k=20 generation ablation below.

\begin{figure}
\centering
\pandocbounded{\includegraphics[keepaspectratio,alt={Retrieval strategies at matched k=10.}]{fig_strategy.png}}
\caption{Retrieval strategies at matched k=10.}
\end{figure}

\textbf{Negative results: BM25, hybrid, and reranking.} BM25 reaches
only 0.153 at k=10, one third of dense retrieval, and it drags every
hybrid mixture below pure dense (best hybrid 0.320 at alpha 0.7).
Filings are boilerplate: the terms in a question (``capital
expenditure'', ``fiscal 2018'') appear on dozens of pages of every
document, so term frequency barely discriminates, while paraphrase (the
statement line reads ``purchases of property, plant and equipment'')
defeats exact matching outright. The ms-marco cross-encoder also
underperforms the dense ranking it reranks (0.380 to 0.447 versus 0.500
at k=10), and enlarging its pool makes it worse, the opposite of the
usual pattern: trained on web passages, it mis-scores dense financial
tables, and a larger pool only gives it more distractors to promote.
Both components would be net-negative additions to this system despite
their standard-recipe status.

\subsection{Generation}\label{generation-1}

\begin{longtable}[]{@{}
  >{\raggedright\arraybackslash}p{(\linewidth - 18\tabcolsep) * \real{0.1000}}
  >{\raggedright\arraybackslash}p{(\linewidth - 18\tabcolsep) * \real{0.1000}}
  >{\raggedright\arraybackslash}p{(\linewidth - 18\tabcolsep) * \real{0.1000}}
  >{\raggedright\arraybackslash}p{(\linewidth - 18\tabcolsep) * \real{0.1000}}
  >{\raggedright\arraybackslash}p{(\linewidth - 18\tabcolsep) * \real{0.1000}}
  >{\raggedright\arraybackslash}p{(\linewidth - 18\tabcolsep) * \real{0.1000}}
  >{\raggedright\arraybackslash}p{(\linewidth - 18\tabcolsep) * \real{0.1000}}
  >{\raggedright\arraybackslash}p{(\linewidth - 18\tabcolsep) * \real{0.1000}}
  >{\raggedright\arraybackslash}p{(\linewidth - 18\tabcolsep) * \real{0.1000}}
  >{\raggedright\arraybackslash}p{(\linewidth - 18\tabcolsep) * \real{0.1000}}@{}}
\caption{generation results. Accuracy is the judge's fully-correct rate;
Numeric is scale-tolerant numeric accuracy; Halluc. is the judge's
hallucination rate.}\tabularnewline
\toprule\noalign{}
\begin{minipage}[b]{\linewidth}\raggedright
Run
\end{minipage} & \begin{minipage}[b]{\linewidth}\raggedright
Context
\end{minipage} & \begin{minipage}[b]{\linewidth}\raggedright
k
\end{minipage} & \begin{minipage}[b]{\linewidth}\raggedright
Accuracy
\end{minipage} & \begin{minipage}[b]{\linewidth}\raggedright
Numeric
\end{minipage} & \begin{minipage}[b]{\linewidth}\raggedright
Abstain
\end{minipage} & \begin{minipage}[b]{\linewidth}\raggedright
Cited
\end{minipage} & \begin{minipage}[b]{\linewidth}\raggedright
Cite valid
\end{minipage} & \begin{minipage}[b]{\linewidth}\raggedright
Halluc.
\end{minipage} & \begin{minipage}[b]{\linewidth}\raggedright
s/q
\end{minipage} \\
\midrule\noalign{}
\endfirsthead
\toprule\noalign{}
\begin{minipage}[b]{\linewidth}\raggedright
Run
\end{minipage} & \begin{minipage}[b]{\linewidth}\raggedright
Context
\end{minipage} & \begin{minipage}[b]{\linewidth}\raggedright
k
\end{minipage} & \begin{minipage}[b]{\linewidth}\raggedright
Accuracy
\end{minipage} & \begin{minipage}[b]{\linewidth}\raggedright
Numeric
\end{minipage} & \begin{minipage}[b]{\linewidth}\raggedright
Abstain
\end{minipage} & \begin{minipage}[b]{\linewidth}\raggedright
Cited
\end{minipage} & \begin{minipage}[b]{\linewidth}\raggedright
Cite valid
\end{minipage} & \begin{minipage}[b]{\linewidth}\raggedright
Halluc.
\end{minipage} & \begin{minipage}[b]{\linewidth}\raggedright
s/q
\end{minipage} \\
\midrule\noalign{}
\endhead
\bottomrule\noalign{}
\endlastfoot
qwen2.5-14B oracle & oracle & 5 & 0.593 & 0.627 & 0.13 & 0.94 & 0.94 &
0.20 & 3.2 \\
qwen2.5-7B oracle & oracle & 5 & 0.427 & 0.437 & 0.34 & 0.63 & 0.63 &
0.32 & 2.7 \\
qwen2.5-14B & retrieved & 20 & 0.387 & 0.421 & 0.19 & 0.60 & 0.27 & 0.28
& 17.6 \\
phi4-14B & retrieved & 10 & 0.367 & 0.381 & 0.49 & 0.66 & 0.31 & 0.28 &
5.2 \\
qwen2.5-14B & retrieved & 10 & 0.367 & 0.381 & 0.36 & 0.65 & 0.34 & 0.27
& 12.3 \\
mistral-7B & retrieved & 10 & 0.253 & 0.222 & 0.51 & 0.47 & 0.17 & 0.33
& 2.8 \\
qwen2.5-7B & retrieved & 10 & 0.247 & 0.262 & 0.55 & 0.43 & 0.21 & 0.31
& 2.3 \\
llama3.1-8B & retrieved & 10 & 0.200 & 0.278 & 0.46 & 0.66 & 0.33 & 0.45
& 3.4 \\
qwen2.5-14B closed book & none & 0 & n/a & 0.111 & 0.00 & 0.00 & 0.00 &
n/a & 12.5 \\
qwen2.5-7B closed book & none & 0 & n/a & 0.087 & 0.00 & 0.00 & 0.00 &
n/a & 2.0 \\
\end{longtable}

\begin{figure}
\centering
\pandocbounded{\includegraphics[keepaspectratio,alt={Generation models under retrieved context.}]{fig_gen_models.png}}
\caption{Generation models under retrieved context.}
\end{figure}

\textbf{Model comparison.} The two 14B models tie at 0.367 under
identical retrieved context, well ahead of the 7B and 8B class (0.200 to
0.253). Within the qwen family, doubling parameters is worth 12 points
(0.247 to 0.367). llama3.1-8B pairs the lowest accuracy with the highest
hallucination rate (0.45) despite citing frequently, the worst
combination for a system meant to be trusted. Notably, phi4-14B matches
qwen2.5-14B while abstaining far more often (0.49 versus 0.36): it
trades coverage for precision.

\begin{figure}
\centering
\pandocbounded{\includegraphics[keepaspectratio,alt={Accuracy under closed book, retrieved, and oracle context.}]{fig_context_modes.png}}
\caption{Accuracy under closed book, retrieved, and oracle context.}
\end{figure}

\textbf{Context ablations.} The closed-book runs collapse to 0.09 to
0.11 numeric accuracy with zero citations, so essentially nothing the
RAG system gets right comes from pre-training memory. The oracle runs
bound what better retrieval could buy: with gold evidence in the prompt,
qwen2.5-14B reaches 0.593 (and cites validly on 94 percent of answers).
Raising k from 10 to 20 lifts context recall from 0.627 to 0.707 but
end-to-end accuracy only from 0.367 to 0.387: the generator recovers
just a fifth of the extra evidence it is given, at a 43 percent latency
premium.

\textbf{Citation behaviour.} Under oracle context, citation validity
equals citation rate (0.94 for the 14B): when the right evidence is
present, the model cites it. Under retrieved context, validity drops to
roughly half the citation rate (0.27 to 0.34 valid against 0.60 to 0.66
cited): models keep citing fluently even when the cited passage does not
contain the evidence, which is exactly the failure mode citation
validity was designed to expose.

\section{Error Analysis}\label{error-analysis}

\subsection{Retrieval failures}\label{retrieval-failures}

At the filtered k=10 operating point, 56 of 150 questions (37.3 percent)
miss the gold page:

{\def\LTcaptype{none} % do not increment counter
\begin{longtable}[]{@{}lll@{}}
\toprule\noalign{}
Failure mode & Count & Share \\
\midrule\noalign{}
\endhead
\bottomrule\noalign{}
\endlastfoot
Wrong document retrieved & 3 & 2.0\% \\
Right document, wrong page & 53 & 35.3\% \\
Gold page hit but evidence text mostly absent & 19 & 12.7\% \\
\end{longtable}
}

Document discrimination is essentially solved by the filter; what
remains is ranking inside a single filing, where the query must
distinguish the income statement from ten pages that mention the same
terms. The third row is invisible to page-level metrics: the right page
was retrieved but the chunk boundary cut the evidence away (median
coverage 0.03 among those 19). Failure rates differ sharply by question
type: 16 percent for metrics-generated questions, whose phrasing mirrors
statement line items, versus 56 percent for domain-relevant and 40
percent for novel-generated questions, whose analyst phrasing shares
little vocabulary with the evidence.

\subsection{Attributing end-to-end
errors}\label{attributing-end-to-end-errors}

Comparing qwen2.5-14B under retrieved (k=10) and oracle context on the
same 150 questions:

{\def\LTcaptype{none} % do not increment counter
\begin{longtable}[]{@{}
  >{\raggedright\arraybackslash}p{(\linewidth - 6\tabcolsep) * \real{0.2500}}
  >{\raggedright\arraybackslash}p{(\linewidth - 6\tabcolsep) * \real{0.2500}}
  >{\raggedright\arraybackslash}p{(\linewidth - 6\tabcolsep) * \real{0.2500}}
  >{\raggedright\arraybackslash}p{(\linewidth - 6\tabcolsep) * \real{0.2500}}@{}}
\toprule\noalign{}
\begin{minipage}[b]{\linewidth}\raggedright
Outcome
\end{minipage} & \begin{minipage}[b]{\linewidth}\raggedright
Count
\end{minipage} & \begin{minipage}[b]{\linewidth}\raggedright
Share
\end{minipage} & \begin{minipage}[b]{\linewidth}\raggedright
Meaning
\end{minipage} \\
\midrule\noalign{}
\endhead
\bottomrule\noalign{}
\endlastfoot
Solved & 43 & 28.7\% & correct with retrieved and with gold evidence \\
Retrieval-limited & 46 & 30.7\% & correct with gold evidence, wrong with
retrieved \\
Generation-limited & 61 & 40.7\% & wrong even with gold evidence \\
Retrieval-rescued & 12 & 8.0\% & wrong with gold evidence, right with
retrieved \\
\end{longtable}
}

The retrieval-limited share (30.7 percent) is the headroom a perfect
retriever could recover; the generation-limited share (40.7 percent) is
the model's own ceiling and is the larger of the two. The small
retrieval-rescued share is consistent with the near-floor closed-book
accuracy: little arrives from memory.

\begin{figure}
\centering
\pandocbounded{\includegraphics[keepaspectratio,alt={Fully-correct rate by reasoning type for the best RAG run.}]{fig_by_reasoning.png}}
\caption{Fully-correct rate by reasoning type for the best RAG run.}
\end{figure}

Generation-limited failures concentrate in multi-step arithmetic: of the
61, 25 are labelled numerical or numerically-grounded logical reasoning,
14 information extraction, and 22 unlabelled (mostly novel-generated
questions that also require computation). Two representative
oracle-context failures: asked for AES Corporation's FY2022 inventory
turnover, the model set up the correct formula, then mis-computed the
average inventory and produced a wrong ratio; asked for Amazon's FY2017
days payable outstanding, it applied the given formula and landed on
94.78 days against a gold 93.86, an arithmetic slip inside an otherwise
correct derivation. These are not retrieval or comprehension errors;
they are calculator errors, and no retrieval improvement can fix them.

\section{Discussion}\label{discussion}

\textbf{What mattered.} One cheap, training-free intervention (metadata
filtering) outweighed every embedding and chunking change we measured,
which inverts the usual tuning priority for corpora of homogeneous,
near-duplicate documents. Embedding choice came second (e5-large-v2),
chunk geometry third and mostly as a coverage-versus-ranking trade-off.
The standard additions, BM25 hybrid and cross-encoder reranking, were
net negative here, a reminder that recipe components carry domain
assumptions: term-frequency discrimination and web-passage relevance
both fail on boilerplate filings.

\textbf{Where the system is bound.} With the best retrieval we built
(0.707 context recall at k=20), end-to-end accuracy is 0.387 while the
oracle bound is 0.593. Closing the entire retrieval gap would therefore
still leave more than 40 percent of questions failing on generation,
dominated by financial arithmetic. The most promising next steps split
accordingly: on the retrieval side, table-aware chunking and rerankers
trained on financial text to attack the within-document ranking problem;
on the generation side, tool-augmented arithmetic (a calculator the
model must route computations through) or larger unquantised models.

\textbf{Limitations.} Ground truth is page-level, so a chunk from the
right page counts as relevant even if tables on that page mislead the
model, and our strict-recall and coverage metrics only approximate
span-level relevance. The judge is qwen2.5-14B, which also appears as a
system under test; a judge may prefer its own outputs, and although its
verdicts agree substantially with the deterministic numeric metric
(kappa 0.607 over 1,008 answers) and the model ranking is unchanged
under the purely deterministic metrics, no human labels were collected.
All models are Q4 quantised, which may understate their unquantised
accuracy, particularly on arithmetic. Runs are deterministic
(temperature 0) single passes, so we report no variance. The metadata
filter depends on the company being named in the question (127 of 150
here); conversational queries with pronouns would need coreference
handling before this gain transfers. Finally, 23 of 84 documents came
from the benchmark's mirror rather than the original URLs; the manifest
records this, but a fully link-rot-proof reproduction depends on that
mirror persisting.

\section{Conclusion}\label{conclusion}

We built and evaluated a complete, fully local RAG pipeline for
FinanceBench: 84 documents and 12,013 pages ingested with full question
coverage, 41 retrieval configurations, five generation models, and
ablations that separate retrieval from generation failures. The best
system (metadata-filtered dense retrieval with bge-base embeddings
feeding qwen2.5-14B at k=20) answers 38.7 percent of the 150 questions
fully correctly by LLM-judge assessment, with a 0.707 context recall and
near-zero reliance on pre-training memory. The experimental program
yielded three transferable lessons: exploit document identity metadata
before tuning anything else on homogeneous corpora; do not assume BM25
hybrids or general-purpose rerankers transfer to financial text; and
measure the oracle bound early, because past a point, and we are past
it, the generator's arithmetic, not the retriever, is the binding
constraint.

\end{document}
